\section{Introduction}
\label{sec:introduction}

Spaceborne synthetic aperture radar (SAR) provides wide-area maritime
observation independent of daylight and through most weather conditions
~\cite{torres2012sentinel1}.
Those properties make Sentinel-1 imagery useful for detecting vessels that do
not have a contemporaneous matched Automatic Identification System (AIS)
report---so-called \emph{dark vessels}. Such detections can support maritime
safety, fisheries monitoring, and investigation of activity potentially
associated with sanctions evasion. The task is nevertheless difficult:
vessels occupy few pixels in a large SAR scene, sea clutter and coastal
structures produce strong responses, and human-verified dark-vessel labels are
expensive. The xView3-SAR benchmark makes this problem concrete with
Sentinel-1 ground-range-detected imagery, georeferenced point annotations, and
an evaluation protocol based on geographic matching~\cite{paolo2022xview3}.

The scarcity of verified labels motivates transfer learning. A representation
learned from a large source corpus may reduce the amount of task-specific
annotation needed to reach a useful detector, but it is not clear which source
is most helpful. Generic ImageNet pretraining supplies broad visual features;
optical remote-sensing pretraining introduces overhead-scene statistics; and
SAR pretraining additionally shares the sensing modality. These alternatives
also interact with architecture. Vision transformers (ViTs) and convolutional
networks encode different inductive biases, and released remote-sensing
checkpoints are not available under one perfectly method-matched pretraining
recipe. A comparison that changes architecture, optimization, detector, and
pretraining source simultaneously therefore cannot isolate the source-domain
question.

We study label efficiency through two architecture-matched tracks, summarized
in Fig.~\ref{fig:study-design}. The ViT track uses ViT-B/16 and the CNN track
uses the similarly sized ConvNeXt-V2-Base. Within each track, four arms differ
only in initialization: random, ImageNet, optical remote sensing, or SAR.
Every arm is fine-tuned end to end through the same point-native detector,
input representation, optimizer, schedule, label subsets, and fixed seed. The
four training subsets are nested at the scene level, containing exactly 12,
28, 56, and 111 scenes. This produces 32 controlled core cells: eight
initializations at four label fractions. Matched roles across the two tracks
then test whether a within-track pattern persists across architecture
families, while preserving explicit caveats about differences in pretraining
history.

\begin{figure}[t]
\centering
\setlength{\fboxsep}{3pt}
\fbox{\parbox[c][1.75cm][c]{0.26\textwidth}{\centering\scriptsize
\textbf{Frozen scenes}\\[2pt]
111 train / 23 dev / 16 test\\
Nested 12, 28, 56, 111\\
50 verified, reserved once}}
\hfill$\Longrightarrow$\hfill
\fbox{\parbox[c][1.75cm][c]{0.28\textwidth}{\centering\scriptsize
\textbf{Two matched tracks}\\[2pt]
ViT-B/16: random, optical, SAR, ImageNet\\
ConvNeXt-V2-B: same roles}}
\hfill$\Longrightarrow$\hfill
\fbox{\parbox[c][1.75cm][c]{0.25\textwidth}{\centering\scriptsize
\textbf{One detector contract}\\[2pt]
Shared input, head, optimizer, seed, FP32, and point scorer}}
\caption{Controlled study design. Initialization is the only variable within
each architecture track. Matched roles are compared descriptively across
tracks; the final human-verified set remains sealed until the cohort and
analysis code are frozen.}
\label{fig:study-design}
\end{figure}


The study addresses three questions:

\begin{enumerate}
  \item How does random, ImageNet, optical remote-sensing, or SAR
    initialization change held-out detection F1 as the number of labeled
    xView3 scenes increases?
  \item Does SAR pretraining improve upon optical pretraining within each
    architecture, especially in the scarce-label regime, and is the direction
    of that contrast consistent across ViT and CNN tracks?
  \item After model selection and thresholds are frozen, how do the controlled
    arms generalize to human-verified dark vessels and near-shore clutter?
\end{enumerate}

Our contribution is the experimental control rather than a new detector or
pretraining method. Specifically, we provide (i) an eight-arm, two-family
label-efficiency design in which initialization is the only within-track
variable; (ii) a shared point-detection and evaluation contract with
scene-disjoint nested subsets and checkpoint-bound operating thresholds; and
(iii) a once-only evaluation on 50 disjoint human-verified scenes, including
dark-vessel recall and near-shore F1. All six pretrained encoders are released
checkpoints; this work performs no source-domain pretraining.

The reportable core matrix is intentionally uniform: all 32 cells run from
scratch on H100 GPUs in strict IEEE FP32 with TensorFloat-32 disabled, one
process per GPU, effective batch size 16, and seed 0. An earlier V100 campaign
helped reveal an evaluation-conversion defect and a checkpoint/threshold
binding defect. Because both defects could affect early stopping and model
selection, its core outputs are retained only as non-reportable diagnostics
and are never mixed with the corrected H100 cohort. Consequently, until the
corrected cohort and once-only evaluation are complete, this manuscript uses
result placeholders and makes no empirical claim about which initialization
wins.

The remainder of the paper reviews relevant transfer and detection work
(Sect.~\ref{sec:related}), defines the controlled arms and common training
recipe (Sect.~\ref{sec:design}), specifies the data and evaluation protocol
(Sect.~\ref{sec:data-evaluation}), and then presents the registered result and
error-analysis views without changing the evaluation contract.
